\section*{ВВЕДЕНИЕ}
\addcontentsline{toc}{section}{ВВЕДЕНИЕ}

В современных условиях значительная часть взаимодействия между заказчиками и исполнителями услуг переносится в цифровую среду. Пользователи всё чаще предпочитают искать специалистов, сравнивать предложения и оформлять запись через web-сервисы, поскольку такой способ позволяет быстрее получить информацию об услуге, стоимости, исполнителе и доступном времени.

Сфера оказания частных, бытовых и цифровых услуг характеризуется большим количеством исполнителей, различием в условиях выполнения работ, стоимости и уровне доверия между участниками. При ручной организации записи исполнитель вынужден самостоятельно публиковать информацию об услугах, вести расписание, обрабатывать обращения, согласовывать время с заказчиками и фиксировать результаты взаимодействия. Такой подход требует дополнительных временных затрат и повышает вероятность организационных ошибок, связанных с пересечением записей, потерей информации или отсутствием актуальных данных о свободном времени.

Для заказчика важна возможность быстро найти подходящую услугу, ознакомиться с описанием предложения, выбрать удобный временной интервал и создать заявку без длительной предварительной переписки. Поэтому актуальной является разработка web-сервиса, который объединяет публикацию услуг, поиск по параметрам, управление расписанием исполнителя, бронирование свободных слотов и сопровождение заявки в едином интерфейсе.

Разрабатываемая программная система предназначена для размещения объявлений об услугах, поиска исполнителей и организации записи заказчиков на свободные временные интервалы. В системе предусматриваются роли неавторизованного пользователя, заказчика и исполнителя. Неавторизованный пользователь может просматривать каталог и страницы услуг. Заказчик получает возможность искать услуги, добавлять их в избранное, создавать бронирования, вести переписку по заявке и оставлять отзывы. Исполнитель может публиковать услуги, управлять расписанием, обрабатывать входящие заявки, изменять их статусы и оценивать взаимодействие с заказчиком.

Цель настоящей работы -- разработка web-сервиса для размещения объявлений об услугах, обеспечивающего публикацию услуг исполнителями, поиск подходящих предложений заказчиками, бронирование свободных временных слотов и сопровождение взаимодействия между пользователями. Для достижения поставленной цели необходимо решить следующие задачи:
\begin{itemize}
	\item провести анализ предметной области;
	\item определить требования к web-сервису с учётом ролей заказчика, исполнителя и неавторизованного пользователя;
	\item спроектировать структуру web-сервиса, базу данных, архитектуру программной системы и пользовательские сценарии;
	\item реализовать программную систему на основе web-технологий с применением базы данных и средств локального развертывания;
	\item разработать функциональные модули регистрации, авторизации, управления услугами, расписанием, бронированиями, сообщениями, избранным и отзывами;
	\item провести тестирование основных функций web-сервиса и проверить корректность пользовательских сценариев.
\end{itemize}

Структура и объем работы. Отчет состоит из введения, 4 разделов основной части, заключения, списка использованных источников и 2 приложений. Текст выпускной квалификационной работы равен \formbytotal{lastpage}{страниц}{е}{ам}{ам}.

Во введении обоснована актуальность разработки, сформулированы цель и задачи работы. 

В первом разделе проводится анализ предметной области, рассматриваются проблемы ручной организации записи на услуги и выполняется сравнительный анализ существующих интернет-сервисов. 

Во втором разделе приводится техническое задание на разработку web-сервиса, включая требования к функциональности, интерфейсу, надежности, информационной безопасности и совместимости. 

В третьем разделе представлены проектные решения: выбранные технологии реализации, архитектура программной системы, структура базы данных и элементы пользовательского интерфейса. 

В четвертом разделе описаны программные модули web-сервиса и выполнено тестирование разработанного программного продукта.
В заключении излагаются основные результаты работы, полученные в ходе разработки web-сервиса для размещения объявлений об услугах.

В приложении А представлен графический материал.
В приложении Б представлены фрагменты исходного кода.